\section{Yang--Mills property-transport table}
\label{app:ymtable}

The table records what each representation naturally displays and the
additional theorem needed for the next physical inference.  ``Manifest'' does
not mean unique to that representation; ``open'' refers to the desired
four-dimensional pure Yang--Mills continuum chain.

\small
\setlength{\tabcolsep}{3pt}
\begin{longtable}{p{.14\textwidth}p{.16\textwidth}p{.18\textwidth}p{.22\textwidth}p{.20\textwidth}}
\toprule
Representation & Manifest structure & Hidden or vulnerable structure & Required recovery theorem & Status for target chain\\
\midrule
\endfirsthead
\toprule
Representation & Manifest structure & Hidden or vulnerable structure & Required recovery theorem & Status for target chain\\
\midrule
\endhead
Lattice links & finite local action; exact finite gauge covariance & orbit identity, global topology, continuum covariance, physical units & gauge descent, chart/sector coverage, observable renormalization & finite statements available; continuum bridge open\\
Gauge quotient & invariant observables and physical redundancy & local sections, simple product coordinates, some boundary data & measurable disintegration, locality recovery, topological completeness & restricted constructions only\\
Blocked/RG variables & scale flow and candidate contraction & microscopic locality, positivity, large fields, boundary charge, history & exact retained-state map, summable reconstruction errors, full-measure polymer/contour control & uniform basin entry open\\
Infinite-volume Euclidean state & thermodynamic correlations & cutoff identity and finite-volume boundary control & tightness, uniqueness or phase selection, volume-uniform estimates & weak-coupling 4D target open\\
Continuum Schwinger functional & Euclidean covariance and distributional correlations & Hamiltonian, particle interpretation, core density & OS axioms, nontriviality, universality, renormalized observable convergence & construction open\\
OS Hilbert space & positive inner product, vacuum, physical semigroup & completeness of chosen gauge-invariant probes & density/state-separation theorem & open for target construction\\
Hamiltonian spectrum & physical energy and mass-gap predicate & relation to regulated rigidity & Theorem~\ref{thm:recovery-spectrum} after all previous certificates & abstract arrow proved; premises open\\
Auxiliary Markov chain & sampling contraction or Poincare inequality & physical time and Hamiltonian interpretation & independent comparison to Euclidean correlations and reconstructed semigroup & no automatic certificate\\
\bottomrule
\end{longtable}
\normalsize

\subsection{Quantifier ledger}

For an estimate intended to feed Corollary~\ref{cor:os}, the preferred order
is:
\[
 \exists m>0\ \exists a_0>0\
 \forall a<a_0\ \forall\text{ admissible sufficiently large }L\
 \forall A,B\in\mathcal D_{a,L}\ \forall t\geq0.
\]
Support sizes and observable norms must be fixed in physical units or
transformed with declared renormalization factors.  The continuum limit must
be taken in a topology strong enough to recover the relevant quadratic forms.
If the proof first fixes $L$ and obtains constants depending exponentially on
$L$, it has not established the displayed statement.

\subsection{False-positive bank}

\begin{enumerate}[label=(\arabic*),itemsep=2pt]
 \item Finite-volume spectral discreteness does not imply an infinite-volume
 gap.
 \item A lattice-unit gap can collapse after division by the lattice spacing.
 \item Decay on an incomplete observable sector can hide low-energy states.
 \item A gap for an auxiliary sampling generator changes under arbitrary time
 rescaling and is not a physical Hamiltonian gap.
 \item Gauge fixing can preserve local calculations while losing positivity
 or global orbit coverage.
 \item A small-field chart does not control the complement without a
 full-measure theorem.
 \item Confinement or an area law is not used as a synonym for a mass gap.
 \item A continuum limit whose nontriviality or regulator identity is unknown
 is not yet the target theory.
\end{enumerate}
